\section{引言与相关工作}
网络流量矩阵描述节点之间的流量需求，缺失项恢复可为离线网络分析提供输入。
在已有完整历史数据、当前窗口只能观测部分流量的条件下，
模型既可利用相关流的时空上下文，也可直接利用本流邻近观测。
前者能够学习跨流关联，后者提供距离明确的局部插值线索。
二者在信息覆盖上并不互斥：时空模型本就可以访问本流观测，
增加一条显式读出路径是否值得，必须通过匹配对照回答。

本文围绕一个问题展开：\emph{本流显式观测读出，能否改善浅层时空模型的精度与执行成本取舍？}
我们采用一层SPIN时空块，另行读取本流前后最近的真实观测，
将两个方向的局部统计量与上下文查询拼接解码，记为SPIN+Direct。
Direct提供的是受约束的读出方式，并不增加原窗口中不存在的信息，
也不保证无损保留观测幅值。
本轮不增加融合门控、新损失或额外记忆，
而是通过分别从头训练的Context-only与Direct-only，检验两条信息路径的使用价值。

时空与局部信息的结合已有基础。
SRMF结合流量的较大范围结构和局部时空约束，并讨论邻近观测补正\cite{zhang2009srmf}；
因此，本文不将全局与局部结合本身作为新概念。
SPIN通过掩蔽时空注意力从稀疏观测构建目标表示\cite{marisca2022spin}，
同时处理本流时间关系与邻流关系。
本文复用其一层编码块，研究额外显式局部读出在这一浅层骨架上的增量，
不把Context-only称为官方SPIN，也不把相对旧四层SPIN的速度变化全部归因于Direct。

同任务的TM-LLM已发表于ISCC 2025，采用对抗式预补全、任务嵌入与提示，
并适配预训练语言模型\cite{jiang2025tmllm}。
其训练依赖和模型规模路线与本文的小型、从头训练模型不同；
这里比较研究定位，不借用其不同协议下的数值建立优劣结论。
2026年预印本Utimac则在局部窗口中估计流量的统计分布与稀疏偏差，
通过块坐标下降补全并报告不确定性区间，主要面向数据中心流量，附录包含G\'EANT\cite{liu2026utimac}。
本文只输出点估计，不提出新的概率模型或不确定性保证。

此前我们还考察过在SPIN+Direct上增加同步矩阵记忆的Full变体。
该分支延续delta-rule残差更新思想\cite{yang2024deltanet}，
现作为历史扩展对照保留，不承担本轮的主要方法主张。
本文通过共同数据、配对初始化、固定优化日程及实际推理测量，
回答局部读出与上下文组合的有限问题。
结果属于开发评价：同一范围参与检查点选择和研究决策，
因而本文仍是一份具有明确证据边界的内部方法初稿。
